\documentclass[10pt,aspectratio=169]{beamer}
\usepackage[utf8]{inputenc}
\usepackage[T1]{fontenc}
\usepackage[french]{babel}
\usepackage{amsmath}
\usepackage{amssymb}
\usepackage{booktabs}
\usepackage{tikz}
\usepackage{array}
\usepackage{multirow}
\usepackage{colortbl}


% Thème
\usetheme{Madrid}
\usecolortheme{default}

% Couleurs personnalisées
\definecolor{bleuactu}{RGB}{0,51,102}
\definecolor{vertactu}{RGB}{0,102,51}
\definecolor{rougeactu}{RGB}{153,0,0}
\definecolor{codegreen}{rgb}{0,0.6,0}
\definecolor{codegray}{rgb}{0.5,0.5,0.5}
\definecolor{codepurple}{rgb}{0.58,0,0.82}
\definecolor{backcolour}{rgb}{0.95,0.95,0.92}

\setbeamercolor{structure}{fg=bleuactu}
\setbeamercolor{frametitle}{bg=bleuactu,fg=white}

% Informations du document
\title{SCR, MCR et Fonds Propres en Assurance}
\subtitle{Guide Complet Solvabilité II}
\author[F. Borel-Mathurin]{Fabrice BOREL-MATHURIN}
\institute[ACPR-DEAR et SAF Lyon 1] {
ACPR - DEAR\\
et\\
Laboratoire SAF Université Claude Bernard Lyon 1}
\date{6 novembre 2025}
%\date{\today}


\begin{document}

\frame{\titlepage}

\begin{frame}
\frametitle{Sommaire}
\tableofcontents
\end{frame}

\section{Le Capital de Solvabilité Requis (SCR)}
\begin{frame}{SCR - Définition et objectif}
\begin{block}{Définition du SCR}
Montant de capital que doit détenir un assureur pour absorber les pertes imprévues sur un horizon d'un an avec un niveau de confiance de \textbf{99,5\%}
\end{block}

\begin{columns}[T]
\begin{column}{0.48\textwidth}
\textbf{Sources réglementaires:}
\begin{itemize}
\item Directive Solvabilité II (2009/138/CE), Article 101
\item Règlement Délégué (UE) 2015/35, Article 24
\end{itemize}
\end{column}
\begin{column}{0.48\textwidth}
\textbf{Concept économique:}
\begin{itemize}
\item Value-at-Risk (VaR) à 99,5\%
\item Variation des fonds propres sur un an
\end{itemize}
\end{column}
\end{columns}
\end{frame}

\begin{frame}{Structure modulaire du SCR}
\begin{block}{Formule standard}
\[
SCR = BSCR + SCR_{\text{opérationnel}} - Adj_{\text{PT}} - Adj_{\text{DT}}
\]
\end{block}

\begin{itemize}
\item \textbf{BSCR}: Basic SCR (agrégation des modules de risque)
\item \textbf{SCR\textsubscript{opérationnel}}: Risque opérationnel
\item \textbf{Adj\textsubscript{PT}}: Ajustement provisions techniques
\item \textbf{Adj\textsubscript{DT}}: Ajustement impôts différés
\end{itemize}
\end{frame}

\begin{frame}{Modules de risque du BSCR}
\begin{block}{Les 6 modules principaux}
\begin{enumerate}
\item Risque de marché (SCR\textsubscript{market})
\item Risque de contrepartie (SCR\textsubscript{counterparty})
\item Risque de souscription Vie (SCR\textsubscript{life})
\item Risque de souscription Non-Vie (SCR\textsubscript{non-life})
\item Risque de souscription Santé (SCR\textsubscript{health})
\item Risque d'immobilisations incorporelles (SCR\textsubscript{intangible})
\end{enumerate}
\end{block}
\end{frame}

\begin{frame}{Matrice de corrélation niveau 1}
\scriptsize
\begin{table}
\begin{tabular}{l|cccccc}
\textbf{Modules} & \textbf{Marché} & \textbf{Contrepartie} & \textbf{Vie} & \textbf{Non-Vie} & \textbf{Santé} & \textbf{Intangibles} \\
\midrule
\textbf{Marché} & 1,00 & 0,25 & 0,25 & 0,25 & 0,25 & 0,25 \\
\textbf{Contrepartie} & 0,25 & 1,00 & 0,25 & 0,50 & 0,25 & 0,25 \\
\textbf{Vie} & 0,25 & 0,25 & 1,00 & 0,25 & 0,25 & 0,25 \\
\textbf{Non-Vie} & 0,25 & 0,50 & 0,25 & 1,00 & 0,25 & 0,25 \\
\textbf{Santé} & 0,25 & 0,25 & 0,25 & 0,25 & 1,00 & 0,25 \\
\textbf{Intangibles} & 0,25 & 0,25 & 0,25 & 0,25 & 0,25 & 1,00 \\
\end{tabular}
\end{table}
\normalsize
\textbf{Référence}: Règlement Délégué (UE) 2015/35, Annexe IV
\end{frame}

\section{Modules de Risque : Calculs Détaillés}
\begin{frame}{Risque de marché - Sous-modules}
\begin{block}{Risque action (SCR\textsubscript{eq})}
\begin{itemize}
\item \textbf{Actions Type 1} (côtées EEE/OCDE): Choc de 39\% + symmetric adjustment
\item \textbf{Actions Type 2} (autres): Choc de 49\% + symmetric adjustment
\item \textbf{Calcul}: $\Delta NAV = (\text{Actifs actions} \times \text{choc}) - (\Delta BE \text{ lié aux actions})$
\end{itemize}
\end{block}

\begin{block}{Risque immobilier (SCR\textsubscript{prop})}
\begin{itemize}
\item Choc de -25\% sur la valeur des actifs immobiliers
\item Impact sur BE si détention d'actifs immobiliers
\end{itemize}
\end{block}
\end{frame}

\begin{frame}{Risque de taux et de spread}
\begin{block}{Risque de taux (SCR\textsubscript{int})}
\begin{itemize}
\item Deux scénarios (hausse et baisse des taux)
\item Courbe de choc EIOPA par maturité
\item Impact sur obligations + effet inverse sur BE
\end{itemize}
\end{block}

\begin{block}{Risque de spread (SCR\textsubscript{spread})}
Formule: $\text{Spread}_{\text{choc}} = a + b \times \text{duration modifiée}$

\begin{table}
\scriptsize
\begin{tabular}{lc}
\textbf{Rating} & \textbf{Coefficients} \\
\midrule
AAA-AA & a = 0,9\%, b = 0,06\% \\
A & a = 1,1\%, b = 0,07\% \\
BBB & a = 1,4\%, b = 0,08\% \\
BB & a = 3,0\%, b = 0,22\% \\
B & a = 6,5\%, b = 0,38\% \\
C & a = 13,5\%, b = 0,50\% \\
\end{tabular}
\end{table}
\end{block}
\end{frame}

\section{Minimum Capital Requirement (MCR)}
\begin{frame}{MCR - Définition et objectif}
\begin{block}{Définition}
Le MCR représente le \textbf{seuil absolu} en dessous duquel l'assureur ne peut plus exercer son activité
\end{block}

\textbf{Référence}: Directive Solvabilité II, Articles 129-131

\begin{alertblock}{Importance critique}
\begin{itemize}
\item Niveau minimum de capital réglementaire
\item Déclenchement de mesures de sauvegarde
\item Risque de retrait d'agrément
\end{itemize}
\end{alertblock}
\end{frame}

\begin{frame}{Calcul du MCR}
\begin{block}{Formules par ligne de business}
\textbf{Assurance Vie:}
\[
MCR_{\text{Vie}} = 0,04 \times BE_{\text{Vie}} + 0,003 \times \text{Capitaux sous risque}
\]

\textbf{Assurance Non-Vie:}
\[
MCR_{\text{Non-Vie}} = 0,08 \times \text{Primes annuelles} + 0,06 \times BE_{\text{Non-Vie}}
\]
\end{block}

\begin{block}{Encadrement réglementaire}
\[
MCR = \max(MCR_{\text{floor}} ; \min(MCR_{\text{combined}} ; MCR_{\text{cap}}))
\]
\begin{itemize}
\item $MCR_{\text{floor}}$: 2,5 à 3,7 M€
\item $MCR_{\text{cap}}$: 45\% du SCR
\item $MCR_{\text{combined}}$: Somme des MCR par ligne
\end{itemize}
\end{block}
\end{frame}

\begin{frame}{Comparaison SCR vs MCR}
\begin{table}
\begin{tabular}{l|p{4cm}|p{4cm}}
\textbf{Aspect} & \textbf{SCR} & \textbf{MCR} \\
\midrule
\textbf{Niveau de confiance} & 99,5\% sur 1 an & Non spécifié \\
\textbf{Objectif} & Seuil d'alerte & Seuil critique \\
\textbf{Conséquences} & Plan de rétablissement & Plan de financement ou retrait d'agrément \\
\textbf{Typique} & 100-200\% du MCR & 25-35\% du SCR \\
\end{tabular}
\end{table}
\end{frame}

\section{Fonds Propres Prudentiels}
\begin{frame}{Classification en tiers}
\begin{block}{Tier 1 - Capital de base}
\begin{itemize}
\item \textbf{Tier 1 non restreint}: Capital social, primes d'émission
\item \textbf{Tier 1 restreint}: Fonds de dotation mutuelles
\item \textbf{Tier 1 auxiliaire}: Capital callable
\end{itemize}
\end{block}

\begin{block}{Tier 2 - Capital complémentaire}
\begin{itemize}
\item Dettes subordonnées (durée ≥ 10 ans)
\item Lettres de crédit, garanties
\end{itemize}
\end{block}

\begin{block}{Tier 3 - Autres éléments}
\begin{itemize}
\item Dettes subordonnées (durée 5-10 ans)
\item Impôts différés actifs nets
\end{itemize}
\end{block}
\end{frame}

\begin{frame}{Limites d'éligibilité}
\begin{columns}[T]
\begin{column}{0.48\textwidth}
\textbf{Pour couvrir le SCR:}
\begin{itemize}
\item Tier 1 ≥ 50\% du SCR
\item Tier 1 restreint + auxiliaire ≤ 20\% du SCR
\item Tier 3 ≤ 15\% du SCR
\item Tier 2 ≤ 50\% du SCR
\end{itemize}
\end{column}
\begin{column}{0.48\textwidth}
\textbf{Pour couvrir le MCR:}
\begin{itemize}
\item Tier 1 non restreint ≥ 80\% du MCR
\item Tier 2 ≤ 20\% du MCR
\item Tier 3 non éligible
\end{itemize}
\end{column}
\end{columns}
\end{frame}

\section{Travaux Dirigés}
\begin{frame}{Exercice: Calcul du SCR marché}
\begin{exampleblock}{Données}
\begin{itemize}
\item Actions Type 1: 150 M€
\item Actions Type 2: 30 M€
\item Obligations souveraines: 300 M€
\item Obligations corporate (rating A, duration 5 ans): 200 M€
\item Immobilier: 80 M€
\item BE = 750 M€ (90\% sensible aux taux)
\item Symmetric adjustment: -3\%
\end{itemize}
\end{exampleblock}
\end{frame}

\begin{frame}[fragile]{Solution Excel}
\begin{block}{SCR Equity}
\begin{lstlisting}[language=Excel]
= (150 * (39% - 3%) + 30 * (49% - 3%)) * (1 - 5%)
= 150*0,36 + 30*0,46 * 0,95 = 68,4 M€
\end{lstlisting}
\end{block}

\begin{block}{SCR Interest Rate}
\begin{itemize}
\item Scénario hausse: Impact obligations = -4,2\%, Impact BE = +3,8\%
\item Scénario baisse: Impact obligations = +4,5\%, Impact BE = -4,1\%
\item SCR\_Int = max(Scénario hausse; Scénario baisse) = 45,2 M€
\end{itemize}
\end{block}
\end{frame}

\begin{frame}[fragile]{Analyse de sensibilité Python}
\begin{lstlisting}[language=Python]
import matplotlib.pyplot as plt
import numpy as np

parametres = ['Taux', 'Actions', 'Spread', 'Mortalité', 'Rachat']
scr_base = [45, 68, 23, 18, 22]
scr_sensibilite = [52, 81, 28, 21, 29]

plt.figure(figsize=(10, 6))
plt.barh(parametres, np.array(scr_sensibilite) - np.array(scr_base))
plt.xlabel('Impact sur le SCR (M€)')
plt.title('Analyse de sensibilité du SCR')
plt.grid(True, alpha=0.3)
plt.show()
\end{lstlisting}
\end{frame}

\section{Conclusion}
\begin{frame}{Sources Réglementaires Principales}
\begin{itemize}
\item \textbf{Directive 2009/138/CE} (Solvency II)
\item \textbf{Règlement Délégué (UE) 2015/35}
\item \textbf{EIOPA} - Guidelines on SCR standard formula (2014)
\item \textbf{ACPR} - Instruction 2015-I-13 sur le calcul du SCR
\item \textbf{EIOPA} - Technical documentation of the SCR standard formula
\end{itemize}
\end{frame}

\begin{frame}
\centering
\Large
\textbf{Merci pour votre attention}
\end{frame}

\end{document}